\section{Los menores de torres al orden fijo que exige el patrón}
Para un corte entero $p\ge2$ fijado, se requiere en los menores físicos, de tamaño $m$,
\[
 R_T=\log\per-L_T-u_T=\sum_{l=2}^{p-1}m^{-l}F_l(z)+O_p(m^{-p}) .
\]
Recogemos la justificación de este uso a orden arbitrario \emph{fijo}.
En la banda difusa de R1 la suma absoluta de los coeficientes bajos para cada grupo de
Wick con $E$ aristas y $b$ índices tiene cota $(CE/m)^{E-b}$ para $E\le D\sqrt m$,
uniformemente en el intervalo de mezcla $I$ en torno a uno. La cota incluye el número
de contracciones del grupo (R1, \emph{permanente}, extracción por número de aristas).
Poniendo $e=E-b$, los grados mínimos tres dan $E\le3e$. Se retienen $e<p$.
Los grupos con $p\le e<2p$ son finitos y cuestan $O_p(m^{-p})$.
En los demás se suma a lo sumo
$\sum_{e\ge2p}(3e+1)(CD/\sqrt m)^e=O_p(m^{-p})$.
La cola de grados del polinomio completo es exponencial en $-\sqrt m$ por R1.
La de los excesos retenidos también: cada uno consta de finitísimas contracciones
en el propagador analítico del mismo disco, que cuesta $O(1/m)$ por entrada.
Esto controla la sustracción de términos, incluso desconectados, antes del logaritmo.

En $I$, cada $d(y)V_e(y)$ y sus derivadas de un orden fijo son $O_p(m^{-e})$ en un
entorno menor. Para $X=m/T$ con $T$ gamma de forma $m+1$ y tasa uno, se integra su
Taylor hasta grado $2(p-e)-1$. El momento absoluto de orden $2r$ de $X-1$ es
$O_r(m^{-r})$. Se obtiene, por ejemplo, de la recurrencia de R2,
$(m-l)\mu_{l+1}=l(2\mu_l+\mu_{l-1})$, $\mu_0=1,\mu_1=0$.
Los momentos $\E X^d=\prod_{i=0}^{d-1}(1-i/m)^{-1}$ dan series en potencias enteras.
Fuera de $I$ se aplica la cota exterior al polinomio completo. Sólo los polinomios
de Taylor se reintegran allí usando Cauchy--Schwarz, momentos dobles y la cola gamma.
Nunca se integra un determinante a través de un polo. El resto físico es $O_p(m^{-p})$.

La fórmula de R9, partes 7--8, se prolonga ahora conservando todos esos coeficientes
finitos, incluido Stirling al mismo orden. El término basal del permanente antes del logaritmo
es $d(1)\ge1$ en los vértices físicos. En la parametrización normalizada, cada $V_e$ es
suma de integrales de núcleos acotados; la inversión sobre particiones que impone
índices distintos aporta $u^{E-|\pi|}$, $u=1/m$, $E-|\pi|\ge e$.
En cada átomo $X_0$ elegible se borra una lista $I_g$ en cada componente. En ambos ejes se
sustraen masas $h_g$ en sus índices y se divide por $1-\sum h_g|I_g|$;
se usa $u=s/(1-\sum h_g|I_g|)$. Los vértices $h_g=0,s$ recuperan esos menores
(incluidos los cruces basales ausentes). El reequilibrio en medidas deformadas usa el
mismo núcleo basal. En un polidisco de radio fijo para $v,p$ los resolventes y los
núcleos quedan acotados como en la prueba de tres componentes, parte 4. La serie de
Fredholm se domina en grado $l$ por $l^{l/2}C^l/l!$ con medidas de variación acotada.
La brecha, el control de la traza y la reducción del radio separan los denominadores de
cero. Así los coeficientes $F_l$ son holomorfos acotados en estas masas; la desaparición
de los órdenes cero y uno en las máscaras físicas es la de R9 por replicación.

El corrector $u_T$ es $u$ veces una función acotada holomorfa en esas mismas medidas.
Tras quitar el valor vacío, a orden $w$ la dependencia de $d_T=\Delta(u_T+R_T)$
involucra a lo sumo $w-1$ componentes por término: cada masa distinta cuesta $s$ y
ya hay al menos un factor $s$ delante. Se toman tantos Taylor como exige el orden fijo.
El resto real se usa en los vértices sin diferenciarlo. La suma finita sobre los átomos
lo mantiene acotado en las actividades complejas locales: el normalizador está cerca de
uno al reducir el radio para $v$ fijo. No se continúa el tamaño entero del tablero.
